\chapter{39只优先池与公司级验证}

本章不是把39只优先池作为一个统一看多名单，而是把它拆成可执行的验证层。优先池来自142只高影响候选的二次筛选；39只都有Q1财务包，但证据质量、研报时效、现金流和当前价格空间差异很大。

正式模型选择桥按当前价格、概率空间、证据质量、研报日期、外部目标价和原始PDF路径分层。优先池代表进入公司级验证，不代表已经获得正式投资评级。

\scriptsize
\setlength{\tabcolsep}{1.5pt}
\begin{longtable}{L{1.05cm}L{1.55cm}L{1.55cm}R{0.85cm}R{0.85cm}R{0.85cm}R{0.85cm}L{2.25cm}}
\toprule
\textbf{代码} & \textbf{公司} & \textbf{行业} & \textbf{现价} & \textbf{目标} & \textbf{空间} & \textbf{Q1现金流} & \textbf{处置} \\
\midrule
\endfirsthead
\toprule
\textbf{代码} & \textbf{公司} & \textbf{行业} & \textbf{现价} & \textbf{目标} & \textbf{空间} & \textbf{Q1现金流} & \textbf{处置} \\
\midrule
\endhead
000042 & 中洲控股 & 房地产 & 8.02元 & 9元 & 6.5\% & 0.5 & 静默吸纳优先 \\
600150 & 中国船舶 & 国防军工 & 34.31元 & 83元 & 142.9\% & 81.5 & 低位盈利优先 \\
002379 & 宏桥控股 & 有色金属 & 18.66元 & 26元 & 38.3\% & 92.4 & 低位盈利优先 \\
601600 & 中国铝业 & 有色金属 & 8.98元 & 15元 & 62.8\% & 108.8 & 低位盈利优先 \\
600346 & 恒力石化 & 石油石化 & 15.24元 & 14元 & -6.8\% & 45.5 & 低位盈利优先 \\
000858 & 五粮液 & 食品饮料 & 76.40元 & 75元 & -1.5\% & -25.4 & 低位盈利优先 \\
002460 & 赣锋锂业 & 有色金属 & 51.50元 & 54元 & 4.2\% & 7.9 & 低位盈利优先 \\
002466 & 天齐锂业 & 有色金属 & 47.43元 & 44元 & -6.4\% & 2.6 & 低位盈利优先 \\
601688 & 华泰证券 & 非银金融 & 20.63元 & 23元 & 9.4\% & 359.1 & 低位盈利优先 \\
601336 & 新华保险 & 非银金融 & 65.55元 & 78元 & 19.7\% & 363.1 & 低位盈利优先 \\
600037 & 歌华有线 & 传媒 & 6.80元 & 6元 & -17.4\% & 0.4 & 低位盈利优先 \\
002738 & 中矿资源 & 有色金属 & 49.33元 & 39元 & -20.4\% & 0.2 & 低位盈利优先 \\
601069 & 西部黄金 & 有色金属 & 22.96元 & 13元 & -44.5\% & 2.3 & 低位盈利优先 \\
000567 & 海德股份 & 非银金融 & 6.17元 & 7元 & 20.8\% & 0.3 & 低位盈利优先 \\
002602 & 世纪华通 & 传媒 & 14.26元 & 17元 & 15.8\% & 19.4 & 低位盈利优先 \\
600026 & 中远海能 & 交通运输 & 15.83元 & 21元 & 34.8\% & 37.7 & 低位盈利优先 \\
002048 & 宁波华翔 & 汽车 & 21.43元 & 28元 & 33.0\% & 1.9 & 低位盈利优先 \\
600688 & 上海石化 & 石油石化 & 2.84元 & 3元 & -8.1\% & -7.9 & 低位盈利优先 \\
600095 & 湘财股份 & 非银金融 & 7.81元 & 4元 & -46.2\% & 21.0 & 低位盈利优先 \\
000426 & 兴业银锡 & 有色金属 & 30.56元 & 24元 & -22.9\% & 12.2 & 低位盈利优先 \\
600489 & 中金黄金 & 有色金属 & 20.12元 & 21元 & 5.2\% & -21.7 & 低位盈利优先 \\
000737 & 北方铜业 & 有色金属 & 12.75元 & 15元 & 20.8\% & -13.4 & 低位盈利优先 \\
002756 & 永兴材料 & 有色金属 & 44.90元 & 42元 & -6.2\% & 3.0 & 低位盈利优先 \\
300014 & 亿纬锂能 & 电力设备 & 52.28元 & 59元 & 12.1\% & -3.7 & 低位盈利优先 \\
601061 & 中信金属 & 商贸零售 & 10.82元 & 17元 & 57.7\% & 31.6 & 低位盈利优先 \\
600739 & 辽宁成大 & 医药生物 & 10.95元 & 43元 & 290.4\% & -1.8 & 低位盈利优先 \\
601360 & 三六零 & 计算机 & 9.00元 & 13元 & 45.6\% & -0.6 & 低位盈利优先 \\
600918 & 中泰证券 & 非银金融 & 5.71元 & 5元 & -6.3\% & 49.2 & 低位盈利优先 \\
000975 & 山金国际 & 有色金属 & 19.34元 & 20元 & 5.0\% & 24.9 & 低位盈利优先 \\
000623 & 吉林敖东 & 医药生物 & 18.66元 & 78元 & 317.0\% & 0.1 & 低位盈利优先 \\
600292 & 电投水电 & 公用事业 & 12.50元 & 7元 & -47.5\% & 5.2 & 低位盈利优先 \\
000987 & 越秀资本 & 非银金融 & 7.77元 & 8元 & 2.5\% & 39.1 & 低位盈利优先 \\
000155 & 川能动力 & 公用事业 & 11.09元 & 10元 & -9.5\% & -3.2 & 低位盈利优先 \\
601377 & 兴业证券 & 非银金融 & 6.21元 & 6元 & -0.8\% & 105.9 & 低位盈利优先 \\
600120 & 浙江东方 & 非银金融 & 4.91元 & 5元 & 2.2\% & 0.6 & 低位盈利优先 \\
000703 & 恒逸石化 & 石油石化 & 14.48元 & 23元 & 61.7\% & -25.7 & 启动后仍有空间 \\
000301 & 东方盛虹 & 石油石化 & 11.61元 & 11元 & -3.7\% & 35.3 & 启动后仍有空间 \\
001339 & 智微智能 & 计算机 & 89.21元 & 65元 & -26.7\% & -5.3 & 启动后仍有空间 \\
002558 & 巨人网络 & 传媒 & 29.56元 & 34元 & 15.0\% & 13.2 & 启动后仍有空间 \\
\bottomrule
\end{longtable}
\normalsize

\section{优先池的阅读方式}

39只优先池的价值不在于把每一只都包装成推荐，而在于把“低位、业绩和资金”三个信号拆开验证。行业阶段只描述公开成交单统计下的资金行为，H1预告只描述公司初步测算，Q1现金流用于检查利润是否已经形成现金回收，当前价格和概率目标则用于判断市场是否已经提前交易了预期。

因此，本章的判断顺序是：先看公司为何进入优先池，再看利润质量和现金转换，随后检查价格位置与估值空间，最后给出升级条件和失效条件。任何一项缺失，都只进入条件性观察，不直接转为正式配置。
\subsection{正式审计模型}
正式审计模型共3只，代表性标的包括宏桥控股、川能动力、恒逸石化。该层平均一年价格位置约37.1\%，Q1经营现金流为正的有1只。
正式审计模型不是空间排序结果，而是当前价格、情景分母、外部锚和原始研报可以同时复核的结果；川能动力虽然进入该层，但因概率目标低于现价，承担的是下行风险对照功能。
\subsection{条件性高空间观察}
条件性高空间观察共12只，代表性标的包括中国船舶、中国铝业、新华保险、海德股份、世纪华通、中远海能、宁波华翔、北方铜业等。该层平均一年价格位置约24.1\%，Q1经营现金流为正的有9只。
这一层的共同特征是模型空间很高，但至少缺一个能把空间转成可审计投资结论的关键环节，通常是当前外部目标价、研报时效或原始证据路径。它们优先等待证据补齐，而不是按模型空间直接追价。
\subsection{条件性安全边际观察}
条件性安全边际观察共9只，代表性标的包括中洲控股、赣锋锂业、华泰证券、中金黄金、亿纬锂能、山金国际、越秀资本、浙江东方等。该层平均一年价格位置约23.2\%，Q1经营现金流为正的有7只。
这一层有一定正向空间，但安全边际有限，市场已经反映一部分业绩改善。只有当H2分母不下修、现金流稳定、价格回撤后承接有效，才有必要从观察升级。
\subsection{估值风险与观察}
估值风险与观察共15只，代表性标的包括恒力石化、五粮液、天齐锂业、歌华有线、中矿资源、西部黄金、上海石化、湘财股份等。该层平均一年价格位置约26.8\%，Q1经营现金流为正的有12只。
这一层的共同问题是概率目标低于当前价、盈利质量不足或价格已经提前反映预期。它们保留在优先池中用于监控行业和盈利变化，但当前不支持正向配置结论。
\section{39只逐票研究正文}

\paragraph{中洲控股（000042）}
进入逻辑：中洲控股属于房地产，本轮被归入静默吸纳优先，行业阶段为静默吸纳。H1归母净利润中值9.9亿元，扣非净利润中值10.6亿元，同比变化约422.4\%；当前价8.02元，位于一年价格区间25.2\%。
基本面与市场验证：Q1经营现金流为正，利润已有一定现金支撑，但仍需观察H2持续性。近20个交易日涨跌幅约-15.3\%，距离一年高点回撤约-23.3\%。因此该标的的核心问题不是“预告是否增长”，而是销售去化、项目结算、资产处置和经营性现金流能否把预告利润转化为可持续的现金收益。
估值与证据：当前模型概率目标8.54元，相对当前价空间6.5\%，证据质量为中，当前无可审计外部目标价；已使用替代估值或明确证据边界。估值恢复：机构预测来源为同花顺F10，2026E/2027E EPS参考为0.93/1.20元；弱来源/历史目标11.84元，明确赋予0\%权重。normalized EPS / PE with project settlement scenarios：熊/基准/牛正常化EPS分别为0.80、1.10、1.50元；对应倍数为6.0、8.0、9.0倍，对应熊值4.80元、基准值8.80元、牛值13.50元；概率价值8.54元，条件性目标8.54元。该标的当前层级为条件性安全边际观察；未满足正式模型的原因是证据质量不足以支持正式模型；最新研报缺失或早于2026-04-01。
催化剂与失效条件：下一阶段重点跟踪销售去化、项目结算、资产处置和经营性现金流。若出现销售与结算不及预期、资产处置收益不可持续或现金流再次恶化，则H2分母、概率目标和当前处置均应下修。

\paragraph{中国船舶（600150）}
进入逻辑：中国船舶属于国防军工，本轮被归入低位盈利优先，行业阶段为资金观察。H1归母净利润中值101.0亿元，扣非净利润中值99.0亿元，同比变化约167.4\%；当前价34.31元，位于一年价格区间28.9\%。
基本面与市场验证：Q1经营现金流为正，利润已有一定现金支撑，但仍需观察H2持续性。近20个交易日涨跌幅约-6.0\%，距离一年高点回撤约-19.1\%。因此该标的的核心问题不是“预告是否增长”，而是订单交付、合同负债、产品结构和应收账款回款能否把预告利润转化为可持续的现金收益。
估值与证据：当前模型概率目标83.33元，相对当前价空间142.9\%，证据质量为中高，最新研报为2026-07-14，原文未披露目标价；已使用替代估值。该标的当前层级为条件性高空间观察；未满足正式模型的原因是没有可审计的外部目标价/合理价值区间；没有对应的本地原始研报证据路径。
催化剂与失效条件：下一阶段重点跟踪订单交付、合同负债、产品结构和应收账款回款。若出现订单确认延后、交付节奏放缓、应收账款继续扩张或估值透支交付，则H2分母、概率目标和当前处置均应下修。

\paragraph{宏桥控股（002379）}
进入逻辑：宏桥控股属于有色金属，本轮被归入低位盈利优先，行业阶段为单日反弹。H1归母净利润中值155.0亿元，扣非净利润中值154.0亿元，同比变化约75.4\%；当前价18.66元，位于一年价格区间24.8\%。
基本面与市场验证：Q1经营现金流为正，利润已有一定现金支撑，但仍需观察H2持续性。近20个交易日涨跌幅约1.1\%，距离一年高点回撤约-45.8\%。因此该标的的核心问题不是“预告是否增长”，而是金属价格、产量、单位成本、库存和经营现金流能否把预告利润转化为可持续的现金收益。
估值与证据：当前模型概率目标25.80元，相对当前价空间38.3\%，证据质量为高，最新研报为2026-05-13，目标价字段为29.00元。该标的当前层级为正式审计模型；未满足正式模型的原因是正式模型：当前价格、情景估值、外部锚和原始研报证据均可复核。
催化剂与失效条件：下一阶段重点跟踪金属价格、产量、单位成本、库存和经营现金流。若出现价格回落、成本曲线抬升、产量不达预期或利润无法转化为现金，则H2分母、概率目标和当前处置均应下修。

\paragraph{中国铝业（601600）}
进入逻辑：中国铝业属于有色金属，本轮被归入低位盈利优先，行业阶段为单日反弹。H1归母净利润中值117.0亿元，扣非净利润中值115.0亿元，同比变化约65.5\%；当前价8.98元，位于一年价格区间23.8\%。
基本面与市场验证：Q1经营现金流为正，利润已有一定现金支撑，但仍需观察H2持续性。近20个交易日涨跌幅约-17.0\%，距离一年高点回撤约-43.6\%。因此该标的的核心问题不是“预告是否增长”，而是金属价格、产量、单位成本、库存和经营现金流能否把预告利润转化为可持续的现金收益。
估值与证据：当前模型概率目标14.62元，相对当前价空间62.8\%，证据质量为中高，最新研报为2026-04-24，原文未披露目标价；已使用替代估值。该标的当前层级为条件性高空间观察；未满足正式模型的原因是没有可审计的外部目标价/合理价值区间；没有对应的本地原始研报证据路径。
催化剂与失效条件：下一阶段重点跟踪金属价格、产量、单位成本、库存和经营现金流。若出现价格回落、成本曲线抬升、产量不达预期或利润无法转化为现金，则H2分母、概率目标和当前处置均应下修。

\paragraph{恒力石化（600346）}
进入逻辑：恒力石化属于石油石化，本轮被归入低位盈利优先，行业阶段为启动确认。H1归母净利润中值72.0亿元，扣非净利润中值53.3亿元，同比变化约136.1\%；当前价15.24元，位于一年价格区间16.6\%。
基本面与市场验证：Q1经营现金流为正，利润已有一定现金支撑，但仍需观察H2持续性。近20个交易日涨跌幅约-15.5\%，距离一年高点回撤约-40.8\%。因此该标的的核心问题不是“预告是否增长”，而是产品价差、原料成本、装置运行率、库存和现金流能否把预告利润转化为可持续的现金收益。
估值与证据：当前模型概率目标14.21元，相对当前价空间-6.8\%，证据质量为中高，最新研报为2026-07-10，原文未披露目标价；已使用替代估值。该标的当前层级为估值风险/观察；未满足正式模型的原因是没有可审计的外部目标价/合理价值区间；没有对应的本地原始研报证据路径；概率目标低于当前价，仅保留为下行风险对照。
催化剂与失效条件：下一阶段重点跟踪产品价差、原料成本、装置运行率、库存和现金流。若出现价差收窄、原料成本上升、装置检修或盈利恢复无法持续，则H2分母、概率目标和当前处置均应下修。

\paragraph{五粮液（000858）}
进入逻辑：五粮液属于食品饮料，本轮被归入低位盈利优先，行业阶段为单日反弹。H1归母净利润中值89.7亿元，扣非净利润中值87.0亿元，同比变化约93.9\%；当前价76.40元，位于一年价格区间4.4\%。
基本面与市场验证：Q1经营现金流为负或接近零，利润兑现必须经过现金转换验证。近20个交易日涨跌幅约-7.8\%，距离一年高点回撤约-42.4\%。因此该标的的核心问题不是“预告是否增长”，而是销量、渠道库存、批价、终端动销和现金回款能否把预告利润转化为可持续的现金收益。
估值与证据：当前模型概率目标75.26元，相对当前价空间-1.5\%，证据质量为中高，最新研报为2026-05-06，原文未披露目标价；已使用替代估值。该标的当前层级为估值风险/观察；未满足正式模型的原因是没有可审计的外部目标价/合理价值区间；没有对应的本地原始研报证据路径；概率目标低于当前价，仅保留为下行风险对照。
催化剂与失效条件：下一阶段重点跟踪销量、渠道库存、批价、终端动销和现金回款。若出现终端动销弱于预期、渠道库存上升、批价承压或现金回款恶化，则H2分母、概率目标和当前处置均应下修。

\paragraph{赣锋锂业（002460）}
进入逻辑：赣锋锂业属于有色金属，本轮被归入低位盈利优先，行业阶段为单日反弹。H1归母净利润中值41.2亿元，扣非净利润中值36.0亿元，同比变化约876.5\%；当前价51.50元，位于一年价格区间33.0\%。
基本面与市场验证：Q1经营现金流为正，利润已有一定现金支撑，但仍需观察H2持续性。近20个交易日涨跌幅约-26.4\%，距离一年高点回撤约-42.6\%。因此该标的的核心问题不是“预告是否增长”，而是金属价格、产量、单位成本、库存和经营现金流能否把预告利润转化为可持续的现金收益。
估值与证据：当前模型概率目标53.68元，相对当前价空间4.2\%，证据质量为中高，最新研报为2026-07-15，原文未披露目标价；已使用替代估值。估值恢复：机构预测来源为东吴证券原始PDF，2026E/2027E EPS参考为4.45/5.35元；没有可审计正权外部目标，采用House替代估值。cycle-adjusted EPS / PE：熊/基准/牛正常化EPS分别为4.00、4.45、5.35元；对应倍数为10.0、12.0、14.0倍，对应熊值40.00元、基准值53.40元、牛值74.90元；概率价值53.68元，条件性目标53.68元。该标的当前层级为条件性安全边际观察；未满足正式模型的原因是没有可审计的外部目标价/合理价值区间；没有对应的本地原始研报证据路径。
催化剂与失效条件：下一阶段重点跟踪金属价格、产量、单位成本、库存和经营现金流。若出现价格回落、成本曲线抬升、产量不达预期或利润无法转化为现金，则H2分母、概率目标和当前处置均应下修。

\paragraph{天齐锂业（002466）}
进入逻辑：天齐锂业属于有色金属，本轮被归入低位盈利优先，行业阶段为单日反弹。H1归母净利润中值35.5亿元，扣非净利润中值35.0亿元，同比变化约4105.6\%；当前价47.43元，位于一年价格区间31.2\%。
基本面与市场验证：Q1经营现金流为正，利润已有一定现金支撑，但仍需观察H2持续性。近20个交易日涨跌幅约-25.7\%，距离一年高点回撤约-42.0\%。因此该标的的核心问题不是“预告是否增长”，而是金属价格、产量、单位成本、库存和经营现金流能否把预告利润转化为可持续的现金收益。
估值与证据：当前模型概率目标44.40元，相对当前价空间-6.4\%，证据质量为中高，最新研报为2026-05-22，原文未披露目标价；已使用替代估值。该标的当前层级为估值风险/观察；未满足正式模型的原因是没有可审计的外部目标价/合理价值区间；没有对应的本地原始研报证据路径；概率目标低于当前价，仅保留为下行风险对照。
催化剂与失效条件：下一阶段重点跟踪金属价格、产量、单位成本、库存和经营现金流。若出现价格回落、成本曲线抬升、产量不达预期或利润无法转化为现金，则H2分母、概率目标和当前处置均应下修。

\paragraph{华泰证券（601688）}
进入逻辑：华泰证券属于非银金融，本轮被归入低位盈利优先，行业阶段为低位无流。H1归母净利润中值115.1亿元，扣非净利润中值115.2亿元，同比变化约52.5\%；当前价20.63元，位于一年价格区间33.9\%。
基本面与市场验证：Q1经营现金流为正，利润已有一定现金支撑，但仍需观察H2持续性。近20个交易日涨跌幅约-1.4\%，距离一年高点回撤约-21.4\%。因此该标的的核心问题不是“预告是否增长”，而是资本市场成交、两融/权益业务、投资收益和净资本约束能否把预告利润转化为可持续的现金收益。
估值与证据：当前模型概率目标22.56元，相对当前价空间9.4\%，证据质量为中高，最新研报为2026-05-09，原文未披露目标价；已使用替代估值。该标的当前层级为条件性安全边际观察；未满足正式模型的原因是没有可审计的外部目标价/合理价值区间；没有对应的本地原始研报证据路径。
催化剂与失效条件：下一阶段重点跟踪资本市场成交、两融/权益业务、投资收益和净资本约束。若出现市场成交回落、投资收益波动、信用风险上升或估值缺乏安全边际，则H2分母、概率目标和当前处置均应下修。

\paragraph{新华保险（601336）}
进入逻辑：新华保险属于非银金融，本轮被归入低位盈利优先，行业阶段为低位无流。H1归母净利润中值222.0亿元，扣非净利润中值222.7亿元，同比变化约50.0\%；当前价65.55元，位于一年价格区间30.7\%。
基本面与市场验证：Q1经营现金流为正，利润已有一定现金支撑，但仍需观察H2持续性。近20个交易日涨跌幅约2.1\%，距离一年高点回撤约-24.1\%。因此该标的的核心问题不是“预告是否增长”，而是资本市场成交、两融/权益业务、投资收益和净资本约束能否把预告利润转化为可持续的现金收益。
估值与证据：当前模型概率目标78.45元，相对当前价空间19.7\%，证据质量为中高，最新研报为2026-07-14，原文未披露目标价；已使用替代估值。该标的当前层级为条件性高空间观察；未满足正式模型的原因是没有可审计的外部目标价/合理价值区间；没有对应的本地原始研报证据路径。
催化剂与失效条件：下一阶段重点跟踪资本市场成交、两融/权益业务、投资收益和净资本约束。若出现市场成交回落、投资收益波动、信用风险上升或估值缺乏安全边际，则H2分母、概率目标和当前处置均应下修。

\paragraph{歌华有线（600037）}
进入逻辑：歌华有线属于传媒，本轮被归入低位盈利优先，行业阶段为资金观察。H1归母净利润中值2.1亿元，扣非净利润中值2.1亿元，同比变化约2297.2\%；当前价6.80元，位于一年价格区间4.7\%。
基本面与市场验证：Q1经营现金流为正，利润已有一定现金支撑，但仍需观察H2持续性。近20个交易日涨跌幅约-4.7\%，距离一年高点回撤约-33.8\%。因此该标的的核心问题不是“预告是否增长”，而是产品上线、流水/付费、内容成本和经营现金流能否把预告利润转化为可持续的现金收益。
估值与证据：当前模型概率目标5.62元，相对当前价空间-17.3\%，证据质量为中，最新研报为2026-03-23，原文未披露目标价；已使用替代估值。估值恢复：机构预测来源为中邮证券原始PDF，2026E/2027E EPS参考为0.03/0.03元；没有可审计正权外部目标，采用House替代估值。normalized EPS / PE with turnaround discount：熊/基准/牛正常化EPS分别为0.02、0.03、0.05元；对应倍数为120.0、180.0、220.0倍，对应熊值2.40元、基准值5.40元、牛值11.00元；概率价值5.62元，条件性目标5.62元。该标的当前层级为估值风险/观察；未满足正式模型的原因是证据质量不足以支持正式模型；最新研报缺失或早于2026-04-01；没有可审计的外部目标价/合理价值区间；没有对应的本地原始研报证据路径；概率目标低于当前价，仅保留为下行风险对照。
催化剂与失效条件：下一阶段重点跟踪产品上线、流水/付费、内容成本和经营现金流。若出现产品或内容兑现延后、用户/付费不及预期或现金流不能跟随利润，则H2分母、概率目标和当前处置均应下修。

\paragraph{中矿资源（002738）}
进入逻辑：中矿资源属于有色金属，本轮被归入低位盈利优先，行业阶段为单日反弹。H1归母净利润中值11.5亿元，扣非净利润中值11.0亿元，同比变化约1190.3\%；当前价49.33元，位于一年价格区间33.5\%。
基本面与市场验证：Q1经营现金流为正，利润已有一定现金支撑，但仍需观察H2持续性。近20个交易日涨跌幅约-15.9\%，距离一年高点回撤约-45.2\%。因此该标的的核心问题不是“预告是否增长”，而是金属价格、产量、单位成本、库存和经营现金流能否把预告利润转化为可持续的现金收益。
估值与证据：当前模型概率目标39.26元，相对当前价空间-20.4\%，证据质量为中高，最新研报为2026-07-14，原文未披露目标价；已使用替代估值。该标的当前层级为估值风险/观察；未满足正式模型的原因是没有可审计的外部目标价/合理价值区间；没有对应的本地原始研报证据路径；概率目标低于当前价，仅保留为下行风险对照。
催化剂与失效条件：下一阶段重点跟踪金属价格、产量、单位成本、库存和经营现金流。若出现价格回落、成本曲线抬升、产量不达预期或利润无法转化为现金，则H2分母、概率目标和当前处置均应下修。

\paragraph{西部黄金（601069）}
进入逻辑：西部黄金属于有色金属，本轮被归入低位盈利优先，行业阶段为单日反弹。H1归母净利润中值5.3亿元，扣非净利润中值5.3亿元，同比变化约306.8\%；当前价22.96元，位于一年价格区间18.0\%。
基本面与市场验证：Q1经营现金流为正，利润已有一定现金支撑，但仍需观察H2持续性。近20个交易日涨跌幅约-11.3\%，距离一年高点回撤约-50.4\%。因此该标的的核心问题不是“预告是否增长”，而是金属价格、产量、单位成本、库存和经营现金流能否把预告利润转化为可持续的现金收益。
估值与证据：当前模型概率目标12.74元，相对当前价空间-44.5\%，证据质量为中，最新研报为2025-10-31，原文未披露目标价；已使用替代估值。该标的当前层级为估值风险/观察；未满足正式模型的原因是证据质量不足以支持正式模型；最新研报缺失或早于2026-04-01；没有可审计的外部目标价/合理价值区间；没有对应的本地原始研报证据路径；概率目标低于当前价，仅保留为下行风险对照。
催化剂与失效条件：下一阶段重点跟踪金属价格、产量、单位成本、库存和经营现金流。若出现价格回落、成本曲线抬升、产量不达预期或利润无法转化为现金，则H2分母、概率目标和当前处置均应下修。

\paragraph{海德股份（000567）}
进入逻辑：海德股份属于非银金融，本轮被归入低位盈利优先，行业阶段为低位无流。H1归母净利润中值14.5亿元，扣非净利润中值14.4亿元，同比变化约886.2\%；当前价6.17元，位于一年价格区间31.8\%。
基本面与市场验证：Q1经营现金流为正，利润已有一定现金支撑，但仍需观察H2持续性。近20个交易日涨跌幅约2.0\%，距离一年高点回撤约-25.4\%。因此该标的的核心问题不是“预告是否增长”，而是资本市场成交、两融/权益业务、投资收益和净资本约束能否把预告利润转化为可持续的现金收益。
估值与证据：当前模型概率目标7.45元，相对当前价空间20.8\%，证据质量为中，最新研报为2024-06-20，原文未披露目标价；已使用替代估值。该标的当前层级为条件性高空间观察；未满足正式模型的原因是证据质量不足以支持正式模型；最新研报缺失或早于2026-04-01；没有可审计的外部目标价/合理价值区间；没有对应的本地原始研报证据路径。
催化剂与失效条件：下一阶段重点跟踪资本市场成交、两融/权益业务、投资收益和净资本约束。若出现市场成交回落、投资收益波动、信用风险上升或估值缺乏安全边际，则H2分母、概率目标和当前处置均应下修。

\paragraph{世纪华通（002602）}
进入逻辑：世纪华通属于传媒，本轮被归入低位盈利优先，行业阶段为资金观察。H1归母净利润中值45.4亿元，扣非净利润中值43.1亿元，同比变化约70.7\%；当前价14.26元，位于一年价格区间22.9\%。
基本面与市场验证：Q1经营现金流为正，利润已有一定现金支撑，但仍需观察H2持续性。近20个交易日涨跌幅约-2.6\%，距离一年高点回撤约-38.0\%。因此该标的的核心问题不是“预告是否增长”，而是产品上线、流水/付费、内容成本和经营现金流能否把预告利润转化为可持续的现金收益。
估值与证据：当前模型概率目标16.51元，相对当前价空间15.8\%，证据质量为中高，最新研报为2026-06-30，原文未披露目标价；已使用替代估值。该标的当前层级为条件性高空间观察；未满足正式模型的原因是没有可审计的外部目标价/合理价值区间；没有对应的本地原始研报证据路径。
催化剂与失效条件：下一阶段重点跟踪产品上线、流水/付费、内容成本和经营现金流。若出现产品或内容兑现延后、用户/付费不及预期或现金流不能跟随利润，则H2分母、概率目标和当前处置均应下修。

\paragraph{中远海能（600026）}
进入逻辑：中远海能属于交通运输，本轮被归入低位盈利优先，行业阶段为低位无流。H1归母净利润中值45.0亿元，扣非净利润中值44.0亿元，同比变化约141.0\%；当前价15.83元，位于一年价格区间32.7\%。
基本面与市场验证：Q1经营现金流为正，利润已有一定现金支撑，但仍需观察H2持续性。近20个交易日涨跌幅约-20.6\%，距离一年高点回撤约-41.4\%。因此该标的的核心问题不是“预告是否增长”，而是运价、运力利用率、燃料成本、资本开支和现金流能否把预告利润转化为可持续的现金收益。
估值与证据：当前模型概率目标21.34元，相对当前价空间34.8\%，证据质量为中高，最新研报为2026-07-12，原文未披露目标价；已使用替代估值。该标的当前层级为条件性高空间观察；未满足正式模型的原因是没有可审计的外部目标价/合理价值区间；没有对应的本地原始研报证据路径。
催化剂与失效条件：下一阶段重点跟踪运价、运力利用率、燃料成本、资本开支和现金流。若出现运价回落、利用率下降、成本上升或周期盈利被错误年化，则H2分母、概率目标和当前处置均应下修。

\paragraph{宁波华翔（002048）}
进入逻辑：宁波华翔属于汽车，本轮被归入低位盈利优先，行业阶段为单日反弹。H1归母净利润中值6.5亿元，扣非净利润中值5.3亿元，同比变化约274.0\%；当前价21.43元，位于一年价格区间18.3\%。
基本面与市场验证：Q1经营现金流为正，利润已有一定现金支撑，但仍需观察H2持续性。近20个交易日涨跌幅约-26.8\%，距离一年高点回撤约-46.8\%。因此该标的的核心问题不是“预告是否增长”，而是销量、单车盈利、产品结构、价格竞争和应收回款能否把预告利润转化为可持续的现金收益。
估值与证据：当前模型概率目标28.50元，相对当前价空间33.0\%，证据质量为中，最新研报为2025-10-30，原文未披露目标价；已使用替代估值。估值恢复：机构预测来源为国金证券原始PDF，2026E/2027E EPS参考为1.98/2.21元；没有可审计正权外部目标，采用House替代估值。normalized EPS / PE：熊/基准/牛正常化EPS分别为1.60、1.98、2.20元；对应倍数为12.0、15.0、18.0倍，对应熊值19.20元、基准值29.64元、牛值39.60元；概率价值28.50元，条件性目标28.50元。该标的当前层级为条件性高空间观察；未满足正式模型的原因是证据质量不足以支持正式模型；最新研报缺失或早于2026-04-01；没有可审计的外部目标价/合理价值区间；没有对应的本地原始研报证据路径。
催化剂与失效条件：下一阶段重点跟踪销量、单车盈利、产品结构、价格竞争和应收回款。若出现销量或单车盈利下滑、价格战加剧、产品结构恶化或现金流转弱，则H2分母、概率目标和当前处置均应下修。

\paragraph{上海石化（600688）}
进入逻辑：上海石化属于石油石化，本轮被归入低位盈利优先，行业阶段为启动确认。H1归母净利润中值2.9亿元，扣非净利润中值3.1亿元，同比变化约163.3\%；当前价2.84元，位于一年价格区间24.8\%。
基本面与市场验证：Q1经营现金流为负或接近零，利润兑现必须经过现金转换验证。近20个交易日涨跌幅约-2.3\%，距离一年高点回撤约-23.7\%。因此该标的的核心问题不是“预告是否增长”，而是产品价差、原料成本、装置运行率、库存和现金流能否把预告利润转化为可持续的现金收益。
估值与证据：当前模型概率目标2.61元，相对当前价空间-8.1\%，证据质量为中，最新研报为2023-03-30，目标价字段为4.06元。估值恢复：机构预测来源为历史原始研报，当前无有效年度共识，2026E/2027E EPS参考为当前无有效年度EPS共识；弱来源/历史目标4.06元，明确赋予0\%权重。PB / cycle-adjusted asset value：Q1 BPS 2.2324元 × 0.90/1.20/1.50倍PB；历史4.06元目标权重为0\%，对应熊值2.01元、基准值2.68元、牛值3.35元；概率价值2.61元，条件性目标2.61元。该标的当前层级为估值风险/观察；未满足正式模型的原因是证据质量不足以支持正式模型；最新研报缺失或早于2026-04-01；概率目标低于当前价，仅保留为下行风险对照。
催化剂与失效条件：下一阶段重点跟踪产品价差、原料成本、装置运行率、库存和现金流。若出现价差收窄、原料成本上升、装置检修或盈利恢复无法持续，则H2分母、概率目标和当前处置均应下修。

\paragraph{湘财股份（600095）}
进入逻辑：湘财股份属于非银金融，本轮被归入低位盈利优先，行业阶段为低位无流。H1归母净利润中值5.4亿元，扣非净利润中值5.2亿元，同比变化约281.3\%；当前价7.81元，位于一年价格区间6.8\%。
基本面与市场验证：Q1经营现金流为正，利润已有一定现金支撑，但仍需观察H2持续性。近20个交易日涨跌幅约-5.4\%，距离一年高点回撤约-43.9\%。因此该标的的核心问题不是“预告是否增长”，而是资本市场成交、两融/权益业务、投资收益和净资本约束能否把预告利润转化为可持续的现金收益。
估值与证据：当前模型概率目标4.20元，相对当前价空间-46.2\%，证据质量为中，最新研报为2026-03-24，原文未披露目标价；已使用替代估值。该标的当前层级为估值风险/观察；未满足正式模型的原因是证据质量不足以支持正式模型；最新研报缺失或早于2026-04-01；没有可审计的外部目标价/合理价值区间；没有对应的本地原始研报证据路径；概率目标低于当前价，仅保留为下行风险对照。
催化剂与失效条件：下一阶段重点跟踪资本市场成交、两融/权益业务、投资收益和净资本约束。若出现市场成交回落、投资收益波动、信用风险上升或估值缺乏安全边际，则H2分母、概率目标和当前处置均应下修。

\paragraph{兴业银锡（000426）}
进入逻辑：兴业银锡属于有色金属，本轮被归入低位盈利优先，行业阶段为单日反弹。H1归母净利润中值22.6亿元，扣非净利润中值18.1亿元，同比变化约183.4\%；当前价30.56元，位于一年价格区间30.1\%。
基本面与市场验证：Q1经营现金流为正，利润已有一定现金支撑，但仍需观察H2持续性。近20个交易日涨跌幅约-20.6\%，距离一年高点回撤约-53.5\%。因此该标的的核心问题不是“预告是否增长”，而是金属价格、产量、单位成本、库存和经营现金流能否把预告利润转化为可持续的现金收益。
估值与证据：当前模型概率目标23.56元，相对当前价空间-22.9\%，证据质量为中高，最新研报为2026-04-24，原文未披露目标价；已使用替代估值。该标的当前层级为估值风险/观察；未满足正式模型的原因是没有可审计的外部目标价/合理价值区间；没有对应的本地原始研报证据路径；概率目标低于当前价，仅保留为下行风险对照。
催化剂与失效条件：下一阶段重点跟踪金属价格、产量、单位成本、库存和经营现金流。若出现价格回落、成本曲线抬升、产量不达预期或利润无法转化为现金，则H2分母、概率目标和当前处置均应下修。

\paragraph{中金黄金（600489）}
进入逻辑：中金黄金属于有色金属，本轮被归入低位盈利优先，行业阶段为单日反弹。H1归母净利润中值43.5亿元，扣非净利润中值43.0亿元，同比变化约61.4\%；当前价20.12元，位于一年价格区间24.6\%。
基本面与市场验证：Q1经营现金流为负或接近零，利润兑现必须经过现金转换验证。近20个交易日涨跌幅约-9.3\%，距离一年高点回撤约-50.1\%。因此该标的的核心问题不是“预告是否增长”，而是金属价格、产量、单位成本、库存和经营现金流能否把预告利润转化为可持续的现金收益。
估值与证据：当前模型概率目标21.16元，相对当前价空间5.2\%，证据质量为中高，最新研报为2026-05-08，原文未披露目标价；已使用替代估值。该标的当前层级为条件性安全边际观察；未满足正式模型的原因是没有可审计的外部目标价/合理价值区间；没有对应的本地原始研报证据路径。
催化剂与失效条件：下一阶段重点跟踪金属价格、产量、单位成本、库存和经营现金流。若出现价格回落、成本曲线抬升、产量不达预期或利润无法转化为现金，则H2分母、概率目标和当前处置均应下修。

\paragraph{北方铜业（000737）}
进入逻辑：北方铜业属于有色金属，本轮被归入低位盈利优先，行业阶段为单日反弹。H1归母净利润中值13.6亿元，扣非净利润中值13.5亿元，同比变化约180.0\%；当前价12.75元，位于一年价格区间19.0\%。
基本面与市场验证：Q1经营现金流为负或接近零，利润兑现必须经过现金转换验证。近20个交易日涨跌幅约-16.9\%，距离一年高点回撤约-39.1\%。因此该标的的核心问题不是“预告是否增长”，而是金属价格、产量、单位成本、库存和经营现金流能否把预告利润转化为可持续的现金收益。
估值与证据：当前模型概率目标15.40元，相对当前价空间20.8\%，证据质量为中低，当前无可审计外部目标价；已使用替代估值或明确证据边界。该标的当前层级为条件性高空间观察；未满足正式模型的原因是证据质量不足以支持正式模型；最新研报缺失或早于2026-04-01；没有可审计的外部目标价/合理价值区间；没有对应的本地原始研报证据路径。
催化剂与失效条件：下一阶段重点跟踪金属价格、产量、单位成本、库存和经营现金流。若出现价格回落、成本曲线抬升、产量不达预期或利润无法转化为现金，则H2分母、概率目标和当前处置均应下修。

\paragraph{永兴材料（002756）}
进入逻辑：永兴材料属于有色金属，本轮被归入低位盈利优先，行业阶段为单日反弹。H1归母净利润中值10.5亿元，扣非净利润中值10.0亿元，同比变化约162.0\%；当前价44.90元，位于一年价格区间27.0\%。
基本面与市场验证：Q1经营现金流为正，利润已有一定现金支撑，但仍需观察H2持续性。近20个交易日涨跌幅约-29.4\%，距离一年高点回撤约-47.6\%。因此该标的的核心问题不是“预告是否增长”，而是金属价格、产量、单位成本、库存和经营现金流能否把预告利润转化为可持续的现金收益。
估值与证据：当前模型概率目标42.10元，相对当前价空间-6.2\%，证据质量为中高，最新研报为2026-07-14，原文未披露目标价；已使用替代估值。该标的当前层级为估值风险/观察；未满足正式模型的原因是没有可审计的外部目标价/合理价值区间；没有对应的本地原始研报证据路径；概率目标低于当前价，仅保留为下行风险对照。
催化剂与失效条件：下一阶段重点跟踪金属价格、产量、单位成本、库存和经营现金流。若出现价格回落、成本曲线抬升、产量不达预期或利润无法转化为现金，则H2分母、概率目标和当前处置均应下修。

\paragraph{亿纬锂能（300014）}
进入逻辑：亿纬锂能属于电力设备，本轮被归入低位盈利优先，行业阶段为单日反弹。H1归母净利润中值32.5亿元，扣非净利润中值25.2亿元，同比变化约102.5\%；当前价52.28元，位于一年价格区间18.7\%。
基本面与市场验证：Q1经营现金流为负或接近零，利润兑现必须经过现金转换验证。近20个交易日涨跌幅约-11.1\%，距离一年高点回撤约-42.4\%。因此该标的的核心问题不是“预告是否增长”，而是出货量、订单、产品价格、产能利用率和营运资金能否把预告利润转化为可持续的现金收益。
估值与证据：当前模型概率目标58.59元，相对当前价空间12.1\%，证据质量为中高，最新研报为2026-06-26，原文未披露目标价；已使用替代估值。该标的当前层级为条件性安全边际观察；未满足正式模型的原因是没有可审计的外部目标价/合理价值区间；没有对应的本地原始研报证据路径。
催化剂与失效条件：下一阶段重点跟踪出货量、订单、产品价格、产能利用率和营运资金。若出现订单转化慢、价格竞争、利用率不足或库存/应收继续占用现金，则H2分母、概率目标和当前处置均应下修。

\paragraph{中信金属（601061）}
进入逻辑：中信金属属于商贸零售，本轮被归入低位盈利优先，行业阶段为单日反弹。H1归母净利润中值26.2亿元，扣非净利润中值26.4亿元，同比变化约81.2\%；当前价10.82元，位于一年价格区间31.0\%。
基本面与市场验证：Q1经营现金流为正，利润已有一定现金支撑，但仍需观察H2持续性。近20个交易日涨跌幅约-12.3\%，距离一年高点回撤约-37.9\%。因此该标的的核心问题不是“预告是否增长”，而是商品价格、贸易量、库存周转、毛利和经营现金流能否把预告利润转化为可持续的现金收益。
估值与证据：当前模型概率目标17.06元，相对当前价空间57.7\%，证据质量为中，最新研报为2025-09-05，目标价字段为11.25元。该标的当前层级为条件性高空间观察；未满足正式模型的原因是证据质量不足以支持正式模型；最新研报缺失或早于2026-04-01。
催化剂与失效条件：下一阶段重点跟踪商品价格、贸易量、库存周转、毛利和经营现金流。若出现贸易量或价差下行、库存周转恶化、毛利承压或现金流弱于利润，则H2分母、概率目标和当前处置均应下修。

\paragraph{辽宁成大（600739）}
进入逻辑：辽宁成大属于医药生物，本轮被归入低位盈利优先，行业阶段为单日反弹。H1归母净利润中值15.8亿元，扣非净利润中值15.8亿元，同比变化约120.3\%；当前价10.95元，位于一年价格区间9.3\%。
基本面与市场验证：Q1经营现金流为负或接近零，利润兑现必须经过现金转换验证。近20个交易日涨跌幅约-1.8\%，距离一年高点回撤约-23.3\%。因此该标的的核心问题不是“预告是否增长”，而是产品放量、集采/价格、研发进度、许可收入和现金流能否把预告利润转化为可持续的现金收益。
估值与证据：当前模型概率目标42.75元，相对当前价空间290.4\%，证据质量为中，最新研报为2017-10-30，目标价字段为22.79元。该标的当前层级为条件性高空间观察；未满足正式模型的原因是证据质量不足以支持正式模型；最新研报缺失或早于2026-04-01。
催化剂与失效条件：下一阶段重点跟踪产品放量、集采/价格、研发进度、许可收入和现金流。若出现产品放量不及预期、价格压力、研发里程碑延后或一次性收益回落，则H2分母、概率目标和当前处置均应下修。

\paragraph{三六零（601360）}
进入逻辑：三六零属于计算机，本轮被归入低位盈利优先，行业阶段为资金观察。H1归母净利润中值2.2亿元，扣非净利润中值1.6亿元，同比变化约178.0\%；当前价9.00元，位于一年价格区间8.6\%。
基本面与市场验证：Q1经营现金流为负或接近零，利润兑现必须经过现金转换验证。近20个交易日涨跌幅约-5.8\%，距离一年高点回撤约-36.7\%。因此该标的的核心问题不是“预告是否增长”，而是合同订单、项目交付、回款、研发投入和利润率能否把预告利润转化为可持续的现金收益。
估值与证据：当前模型概率目标13.10元，相对当前价空间45.6\%，证据质量为中，最新研报为2023-04-25，原文未披露目标价；已使用替代估值。估值恢复：公开分析师聚合目标均值14.44元，3位分析师区间13.70-15.17元；采用2026E收入95.91亿元的PS情景7.5/10.0/11.0倍，对应10.28/13.70/15.07元，概率价值12.95元，叠加公开共识10\%权重后的条件性目标13.10元。该标的当前层级为条件性高空间观察；未满足正式模型的原因是证据质量不足以支持正式模型；最新研报缺失或早于2026-04-01。
催化剂与失效条件：下一阶段重点跟踪合同订单、项目交付、回款、研发投入和利润率。若出现订单确认延后、回款周期拉长、研发费用率上升或估值继续压缩，则H2分母、概率目标和当前处置均应下修。

\paragraph{中泰证券（600918）}
进入逻辑：中泰证券属于非银金融，本轮被归入低位盈利优先，行业阶段为低位无流。H1归母净利润中值17.5亿元，扣非净利润中值17.0亿元，同比变化约146.3\%；当前价5.71元，位于一年价格区间26.1\%。
基本面与市场验证：Q1经营现金流为正，利润已有一定现金支撑，但仍需观察H2持续性。近20个交易日涨跌幅约5.0\%，距离一年高点回撤约-22.9\%。因此该标的的核心问题不是“预告是否增长”，而是资本市场成交、两融/权益业务、投资收益和净资本约束能否把预告利润转化为可持续的现金收益。
估值与证据：当前模型概率目标5.35元，相对当前价空间-6.3\%，证据质量为中，最新研报为2025-10-31，原文未披露目标价；已使用替代估值。该标的当前层级为估值风险/观察；未满足正式模型的原因是证据质量不足以支持正式模型；最新研报缺失或早于2026-04-01；没有可审计的外部目标价/合理价值区间；没有对应的本地原始研报证据路径；概率目标低于当前价，仅保留为下行风险对照。
催化剂与失效条件：下一阶段重点跟踪资本市场成交、两融/权益业务、投资收益和净资本约束。若出现市场成交回落、投资收益波动、信用风险上升或估值缺乏安全边际，则H2分母、概率目标和当前处置均应下修。

\paragraph{山金国际（000975）}
进入逻辑：山金国际属于有色金属，本轮被归入低位盈利优先，行业阶段为单日反弹。H1归母净利润中值24.1亿元，扣非净利润中值23.4亿元，同比变化约51.3\%；当前价19.34元，位于一年价格区间14.6\%。
基本面与市场验证：Q1经营现金流为正，利润已有一定现金支撑，但仍需观察H2持续性。近20个交易日涨跌幅约-10.2\%，距离一年高点回撤约-48.0\%。因此该标的的核心问题不是“预告是否增长”，而是金属价格、产量、单位成本、库存和经营现金流能否把预告利润转化为可持续的现金收益。
估值与证据：当前模型概率目标20.31元，相对当前价空间5.0\%，证据质量为中高，最新研报为2026-05-13，原文未披露目标价；已使用替代估值。该标的当前层级为条件性安全边际观察；未满足正式模型的原因是没有可审计的外部目标价/合理价值区间；没有对应的本地原始研报证据路径。
催化剂与失效条件：下一阶段重点跟踪金属价格、产量、单位成本、库存和经营现金流。若出现价格回落、成本曲线抬升、产量不达预期或利润无法转化为现金，则H2分母、概率目标和当前处置均应下修。

\paragraph{吉林敖东（000623）}
进入逻辑：吉林敖东属于医药生物，本轮被归入低位盈利优先，行业阶段为单日反弹。H1归母净利润中值20.5亿元，扣非净利润中值22.2亿元，同比变化约60.0\%；当前价18.66元，位于一年价格区间32.3\%。
基本面与市场验证：Q1经营现金流为正，利润已有一定现金支撑，但仍需观察H2持续性。近20个交易日涨跌幅约2.4\%，距离一年高点回撤约-16.9\%。因此该标的的核心问题不是“预告是否增长”，而是产品放量、集采/价格、研发进度、许可收入和现金流能否把预告利润转化为可持续的现金收益。
估值与证据：当前模型概率目标77.82元，相对当前价空间317.0\%，证据质量为中，最新研报为2024-09-03，原文未披露目标价；已使用替代估值。该标的当前层级为条件性高空间观察；未满足正式模型的原因是证据质量不足以支持正式模型；最新研报缺失或早于2026-04-01；没有可审计的外部目标价/合理价值区间；没有对应的本地原始研报证据路径。
催化剂与失效条件：下一阶段重点跟踪产品放量、集采/价格、研发进度、许可收入和现金流。若出现产品放量不及预期、价格压力、研发里程碑延后或一次性收益回落，则H2分母、概率目标和当前处置均应下修。

\paragraph{电投水电（600292）}
进入逻辑：电投水电属于公用事业，本轮被归入低位盈利优先，行业阶段为低位无流。H1归母净利润中值11.9亿元，扣非净利润中值11.8亿元，同比变化约122.0\%；当前价12.50元，位于一年价格区间14.9\%。
基本面与市场验证：Q1经营现金流为正，利润已有一定现金支撑，但仍需观察H2持续性。近20个交易日涨跌幅约-15.8\%，距离一年高点回撤约-26.5\%。因此该标的的核心问题不是“预告是否增长”，而是利用小时、上网电价、燃料/资源成本、资本开支和现金流能否把预告利润转化为可持续的现金收益。
估值与证据：当前模型概率目标6.56元，相对当前价空间-47.5\%，证据质量为中，最新研报为2025-09-12，原文未披露目标价；已使用替代估值。该标的当前层级为估值风险/观察；未满足正式模型的原因是证据质量不足以支持正式模型；最新研报缺失或早于2026-04-01；没有可审计的外部目标价/合理价值区间；没有对应的本地原始研报证据路径；概率目标低于当前价，仅保留为下行风险对照。
催化剂与失效条件：下一阶段重点跟踪利用小时、上网电价、燃料/资源成本、资本开支和现金流。若出现利用率或电价下行、成本抬升、资本开支超预期或现金流不达标，则H2分母、概率目标和当前处置均应下修。

\paragraph{越秀资本（000987）}
进入逻辑：越秀资本属于非银金融，本轮被归入低位盈利优先，行业阶段为低位无流。H1归母净利润中值28.8亿元，扣非净利润中值23.3亿元，同比变化约85.0\%；当前价7.77元，位于一年价格区间22.9\%。
基本面与市场验证：Q1经营现金流为正，利润已有一定现金支撑，但仍需观察H2持续性。近20个交易日涨跌幅约-0.7\%，距离一年高点回撤约-34.5\%。因此该标的的核心问题不是“预告是否增长”，而是资本市场成交、两融/权益业务、投资收益和净资本约束能否把预告利润转化为可持续的现金收益。
估值与证据：当前模型概率目标7.96元，相对当前价空间2.5\%，证据质量为中，最新研报为2019-08-19，目标价字段为9.44元。该标的当前层级为条件性安全边际观察；未满足正式模型的原因是证据质量不足以支持正式模型；最新研报缺失或早于2026-04-01。
催化剂与失效条件：下一阶段重点跟踪资本市场成交、两融/权益业务、投资收益和净资本约束。若出现市场成交回落、投资收益波动、信用风险上升或估值缺乏安全边际，则H2分母、概率目标和当前处置均应下修。

\paragraph{川能动力（000155）}
进入逻辑：川能动力属于公用事业，本轮被归入低位盈利优先，行业阶段为低位无流。H1归母净利润中值6.8亿元，扣非净利润中值6.8亿元，同比变化约123.8\%；当前价11.09元，位于一年价格区间15.9\%。
基本面与市场验证：Q1经营现金流为负或接近零，利润兑现必须经过现金转换验证。近20个交易日涨跌幅约-16.9\%，距离一年高点回撤约-42.0\%。因此该标的的核心问题不是“预告是否增长”，而是利用小时、上网电价、燃料/资源成本、资本开支和现金流能否把预告利润转化为可持续的现金收益。
估值与证据：当前模型概率目标10.04元，相对当前价空间-9.5\%，证据质量为高，最新研报为2026-04-20，目标价字段为20.43元。该标的当前层级为正式审计模型；未满足正式模型的原因是正式下行风险对照：证据链完整，但概率目标低于当前价。
催化剂与失效条件：下一阶段重点跟踪利用小时、上网电价、燃料/资源成本、资本开支和现金流。若出现利用率或电价下行、成本抬升、资本开支超预期或现金流不达标，则H2分母、概率目标和当前处置均应下修。

\paragraph{兴业证券（601377）}
进入逻辑：兴业证券属于非银金融，本轮被归入低位盈利优先，行业阶段为低位无流。H1归母净利润中值21.6亿元，扣非净利润中值20.9亿元，同比变化约62.5\%；当前价6.21元，位于一年价格区间22.3\%。
基本面与市场验证：Q1经营现金流为正，利润已有一定现金支撑，但仍需观察H2持续性。近20个交易日涨跌幅约0.3\%，距离一年高点回撤约-21.1\%。因此该标的的核心问题不是“预告是否增长”，而是资本市场成交、两融/权益业务、投资收益和净资本约束能否把预告利润转化为可持续的现金收益。
估值与证据：当前模型概率目标6.16元，相对当前价空间-0.8\%，证据质量为中高，最新研报为2026-05-05，原文未披露目标价；已使用替代估值。该标的当前层级为估值风险/观察；未满足正式模型的原因是没有可审计的外部目标价/合理价值区间；没有对应的本地原始研报证据路径；概率目标低于当前价，仅保留为下行风险对照。
催化剂与失效条件：下一阶段重点跟踪资本市场成交、两融/权益业务、投资收益和净资本约束。若出现市场成交回落、投资收益波动、信用风险上升或估值缺乏安全边际，则H2分母、概率目标和当前处置均应下修。

\paragraph{浙江东方（600120）}
进入逻辑：浙江东方属于非银金融，本轮被归入低位盈利优先，行业阶段为低位无流。H1归母净利润中值8.7亿元，扣非净利润中值8.4亿元，同比变化约114.2\%；当前价4.91元，位于一年价格区间20.9\%。
基本面与市场验证：Q1经营现金流为正，利润已有一定现金支撑，但仍需观察H2持续性。近20个交易日涨跌幅约8.6\%，距离一年高点回撤约-32.9\%。因此该标的的核心问题不是“预告是否增长”，而是资本市场成交、两融/权益业务、投资收益和净资本约束能否把预告利润转化为可持续的现金收益。
估值与证据：当前模型概率目标5.02元，相对当前价空间2.2\%，证据质量为中低，当前无可审计外部目标价；已使用替代估值或明确证据边界。该标的当前层级为条件性安全边际观察；未满足正式模型的原因是证据质量不足以支持正式模型；最新研报缺失或早于2026-04-01；没有可审计的外部目标价/合理价值区间；没有对应的本地原始研报证据路径。
催化剂与失效条件：下一阶段重点跟踪资本市场成交、两融/权益业务、投资收益和净资本约束。若出现市场成交回落、投资收益波动、信用风险上升或估值缺乏安全边际，则H2分母、概率目标和当前处置均应下修。

\paragraph{恒逸石化（000703）}
进入逻辑：恒逸石化属于石油石化，本轮被归入启动后仍有空间，行业阶段为启动确认。H1归母净利润中值57.5亿元，扣非净利润中值57.2亿元，同比变化约2436.6\%；当前价14.48元，位于一年价格区间70.7\%。
基本面与市场验证：Q1经营现金流为负或接近零，利润兑现必须经过现金转换验证。近20个交易日涨跌幅约12.6\%，距离一年高点回撤约-20.0\%。因此该标的的核心问题不是“预告是否增长”，而是产品价差、原料成本、装置运行率、库存和现金流能否把预告利润转化为可持续的现金收益。
估值与证据：当前模型概率目标23.42元，相对当前价空间61.7\%，证据质量为高，最新研报为2026-07-05，目标价字段为17.05元。该标的当前层级为正式审计模型；未满足正式模型的原因是正式模型：当前价格、情景估值、外部锚和原始研报证据均可复核。
催化剂与失效条件：下一阶段重点跟踪产品价差、原料成本、装置运行率、库存和现金流。若出现价差收窄、原料成本上升、装置检修或盈利恢复无法持续，则H2分母、概率目标和当前处置均应下修。

\paragraph{东方盛虹（000301）}
进入逻辑：东方盛虹属于石油石化，本轮被归入启动后仍有空间，行业阶段为启动确认。H1归母净利润中值46.0亿元，扣非净利润中值44.2亿元，同比变化约1091.0\%；当前价11.61元，位于一年价格区间62.3\%。
基本面与市场验证：Q1经营现金流为正，利润已有一定现金支撑，但仍需观察H2持续性。近20个交易日涨跌幅约-4.2\%，距离一年高点回撤约-16.1\%。因此该标的的核心问题不是“预告是否增长”，而是产品价差、原料成本、装置运行率、库存和现金流能否把预告利润转化为可持续的现金收益。
估值与证据：当前模型概率目标11.18元，相对当前价空间-3.7\%，证据质量为中高，最新研报为2026-05-06，原文未披露目标价；已使用替代估值。该标的当前层级为估值风险/观察；未满足正式模型的原因是没有可审计的外部目标价/合理价值区间；没有对应的本地原始研报证据路径；概率目标低于当前价，仅保留为下行风险对照。
催化剂与失效条件：下一阶段重点跟踪产品价差、原料成本、装置运行率、库存和现金流。若出现价差收窄、原料成本上升、装置检修或盈利恢复无法持续，则H2分母、概率目标和当前处置均应下修。

\paragraph{智微智能（001339）}
进入逻辑：智微智能属于计算机，本轮被归入启动后仍有空间，行业阶段为资金观察。H1归母净利润中值3.8亿元，扣非净利润中值3.7亿元，同比变化约277.1\%；当前价89.21元，位于一年价格区间78.8\%。
基本面与市场验证：Q1经营现金流为负或接近零，利润兑现必须经过现金转换验证。近20个交易日涨跌幅约16.6\%，距离一年高点回撤约-13.9\%。因此该标的的核心问题不是“预告是否增长”，而是合同订单、项目交付、回款、研发投入和利润率能否把预告利润转化为可持续的现金收益。
估值与证据：当前模型概率目标65.43元，相对当前价空间-26.7\%，证据质量为中高，最新研报为2026-06-18，原文未披露目标价；已使用替代估值。该标的当前层级为估值风险/观察；未满足正式模型的原因是没有可审计的外部目标价/合理价值区间；没有对应的本地原始研报证据路径；概率目标低于当前价，仅保留为下行风险对照。
催化剂与失效条件：下一阶段重点跟踪合同订单、项目交付、回款、研发投入和利润率。若出现订单确认延后、回款周期拉长、研发费用率上升或估值继续压缩，则H2分母、概率目标和当前处置均应下修。

\paragraph{巨人网络（002558）}
进入逻辑：巨人网络属于传媒，本轮被归入启动后仍有空间，行业阶段为资金观察。H1归母净利润中值21.0亿元，扣非净利润中值23.0亿元，同比变化约170.2\%；当前价29.56元，位于一年价格区间14.7\%。
基本面与市场验证：Q1经营现金流为正，利润已有一定现金支撑，但仍需观察H2持续性。近20个交易日涨跌幅约10.5\%，距离一年高点回撤约-47.3\%。因此该标的的核心问题不是“预告是否增长”，而是产品上线、流水/付费、内容成本和经营现金流能否把预告利润转化为可持续的现金收益。
估值与证据：当前模型概率目标33.99元，相对当前价空间15.0\%，证据质量为中高，最新研报为2026-07-15，原文未披露目标价；已使用替代估值。该标的当前层级为条件性安全边际观察；未满足正式模型的原因是没有可审计的外部目标价/合理价值区间；没有对应的本地原始研报证据路径。
催化剂与失效条件：下一阶段重点跟踪产品上线、流水/付费、内容成本和经营现金流。若出现产品或内容兑现延后、用户/付费不及预期或现金流不能跟随利润，则H2分母、概率目标和当前处置均应下修。

\section{逐票验证要点摘要}

\begin{longtable}{L{1.15cm}L{1.55cm}L{2.35cm}L{2.55cm}L{2.55cm}}
\toprule
\textbf{代码} & \textbf{标的} & \textbf{进入理由} & \textbf{证据/估值边界} & \textbf{下一季度验证} \\
\midrule
\endfirsthead
\toprule
\textbf{代码} & \textbf{标的} & \textbf{进入理由} & \textbf{证据/估值边界} & \textbf{下一季度验证} \\
\midrule
\endhead
000042 & 中洲控股 & 静默吸纳优先；房地产 / 静默吸纳 & H1扣非 10.6亿；Q1 OCF 0.5亿；证据中；概率空间 6.5\%；normalized EPS / PE with project settlement scenarios：熊/基准/牛正常化EPS分别为0.80、1.10、1.50元；对应倍数为6.0、8.0、9.0倍，对应熊值4.80元、基准值8.80元、牛值13.50元；概率价值8.54元，条件性目标8.54元。 & 需要安全边际；验证H2分母不下修、价格承接有效。 \\
600150 & 中国船舶 & 低位盈利优先；国防军工 / 资金观察 & H1扣非 99.0亿；Q1 OCF 81.5亿；证据中高；概率空间 142.9\% & 补齐当期原始研报锚；验证H2扣非兑现和现金流。 \\
002379 & 宏桥控股 & 低位盈利优先；有色金属 / 单日反弹 & H1扣非 154.0亿；Q1 OCF 92.4亿；证据高；概率空间 38.3\% & 正式模型：H2扣非、现金流与外部锚同时复核。 \\
601600 & 中国铝业 & 低位盈利优先；有色金属 / 单日反弹 & H1扣非 115.0亿；Q1 OCF 108.8亿；证据中高；概率空间 62.8\% & 补齐当期原始研报锚；验证H2扣非兑现和现金流。 \\
600346 & 恒力石化 & 低位盈利优先；石油石化 / 启动确认 & H1扣非 53.3亿；Q1 OCF 45.5亿；证据中高；概率空间 -6.8\% & 当前不升级；先验证盈利质量、估值消化和现金流反转。 \\
000858 & 五粮液 & 低位盈利优先；食品饮料 / 单日反弹 & H1扣非 87.0亿；Q1 OCF -25.4亿；证据中高；概率空间 -1.5\% & 当前不升级；先验证盈利质量、估值消化和现金流反转。 \\
002460 & 赣锋锂业 & 低位盈利优先；有色金属 / 单日反弹 & H1扣非 36.0亿；Q1 OCF 7.9亿；证据中高；概率空间 4.2\%；cycle-adjusted EPS / PE：熊/基准/牛正常化EPS分别为4.00、4.45、5.35元；对应倍数为10.0、12.0、14.0倍，对应熊值40.00元、基准值53.40元、牛值74.90元；概率价值53.68元，条件性目标53.68元。 & 需要安全边际；验证H2分母不下修、价格承接有效。 \\
002466 & 天齐锂业 & 低位盈利优先；有色金属 / 单日反弹 & H1扣非 35.0亿；Q1 OCF 2.6亿；证据中高；概率空间 -6.4\% & 当前不升级；先验证盈利质量、估值消化和现金流反转。 \\
601688 & 华泰证券 & 低位盈利优先；非银金融 / 低位无流 & H1扣非 115.2亿；Q1 OCF 359.1亿；证据中高；概率空间 9.4\% & 需要安全边际；验证H2分母不下修、价格承接有效。 \\
601336 & 新华保险 & 低位盈利优先；非银金融 / 低位无流 & H1扣非 222.7亿；Q1 OCF 363.1亿；证据中高；概率空间 19.7\% & 补齐当期原始研报锚；验证H2扣非兑现和现金流。 \\
600037 & 歌华有线 & 低位盈利优先；传媒 / 资金观察 & H1扣非 2.1亿；Q1 OCF 0.4亿；证据中；概率空间 -17.4\%；normalized EPS / PE with turnaround discount：熊/基准/牛正常化EPS分别为0.02、0.03、0.05元；对应倍数为120.0、180.0、220.0倍，对应熊值2.40元、基准值5.40元、牛值11.00元；概率价值5.62元，条件性目标5.62元。 & 当前不升级；先验证盈利质量、估值消化和现金流反转。 \\
002738 & 中矿资源 & 低位盈利优先；有色金属 / 单日反弹 & H1扣非 11.0亿；Q1 OCF 0.2亿；证据中高；概率空间 -20.4\% & 当前不升级；先验证盈利质量、估值消化和现金流反转。 \\
601069 & 西部黄金 & 低位盈利优先；有色金属 / 单日反弹 & H1扣非 5.3亿；Q1 OCF 2.3亿；证据中；概率空间 -44.5\% & 当前不升级；先验证盈利质量、估值消化和现金流反转。 \\
000567 & 海德股份 & 低位盈利优先；非银金融 / 低位无流 & H1扣非 14.4亿；Q1 OCF 0.3亿；证据中；概率空间 20.8\% & 补齐当期原始研报锚；验证H2扣非兑现和现金流。 \\
002602 & 世纪华通 & 低位盈利优先；传媒 / 资金观察 & H1扣非 43.1亿；Q1 OCF 19.4亿；证据中高；概率空间 15.8\% & 补齐当期原始研报锚；验证H2扣非兑现和现金流。 \\
600026 & 中远海能 & 低位盈利优先；交通运输 / 低位无流 & H1扣非 44.0亿；Q1 OCF 37.7亿；证据中高；概率空间 34.8\% & 补齐当期原始研报锚；验证H2扣非兑现和现金流。 \\
002048 & 宁波华翔 & 低位盈利优先；汽车 / 单日反弹 & H1扣非 5.3亿；Q1 OCF 1.9亿；证据中；概率空间 33.0\%；normalized EPS / PE：熊/基准/牛正常化EPS分别为1.60、1.98、2.20元；对应倍数为12.0、15.0、18.0倍，对应熊值19.20元、基准值29.64元、牛值39.60元；概率价值28.50元，条件性目标28.50元。 & 补齐当期原始研报锚；验证H2扣非兑现和现金流。 \\
600688 & 上海石化 & 低位盈利优先；石油石化 / 启动确认 & H1扣非 3.1亿；Q1 OCF -7.9亿；证据中；概率空间 -8.1\%；PB / cycle-adjusted asset value：Q1 BPS 2.2324元 × 0.90/1.20/1.50倍PB；历史4.06元目标权重为0\%，对应熊值2.01元、基准值2.68元、牛值3.35元；概率价值2.61元，条件性目标2.61元。 & 当前不升级；先验证盈利质量、估值消化和现金流反转。 \\
600095 & 湘财股份 & 低位盈利优先；非银金融 / 低位无流 & H1扣非 5.2亿；Q1 OCF 21.0亿；证据中；概率空间 -46.2\% & 当前不升级；先验证盈利质量、估值消化和现金流反转。 \\
000426 & 兴业银锡 & 低位盈利优先；有色金属 / 单日反弹 & H1扣非 18.1亿；Q1 OCF 12.2亿；证据中高；概率空间 -22.9\% & 当前不升级；先验证盈利质量、估值消化和现金流反转。 \\
600489 & 中金黄金 & 低位盈利优先；有色金属 / 单日反弹 & H1扣非 43.0亿；Q1 OCF -21.7亿；证据中高；概率空间 5.2\% & 需要安全边际；验证H2分母不下修、价格承接有效。 \\
000737 & 北方铜业 & 低位盈利优先；有色金属 / 单日反弹 & H1扣非 13.5亿；Q1 OCF -13.4亿；证据中低；概率空间 20.8\% & 补齐当期原始研报锚；验证H2扣非兑现和现金流。 \\
002756 & 永兴材料 & 低位盈利优先；有色金属 / 单日反弹 & H1扣非 10.0亿；Q1 OCF 3.0亿；证据中高；概率空间 -6.2\% & 当前不升级；先验证盈利质量、估值消化和现金流反转。 \\
300014 & 亿纬锂能 & 低位盈利优先；电力设备 / 单日反弹 & H1扣非 25.2亿；Q1 OCF -3.7亿；证据中高；概率空间 12.1\% & 需要安全边际；验证H2分母不下修、价格承接有效。 \\
601061 & 中信金属 & 低位盈利优先；商贸零售 / 单日反弹 & H1扣非 26.4亿；Q1 OCF 31.6亿；证据中；概率空间 57.7\% & 补齐当期原始研报锚；验证H2扣非兑现和现金流。 \\
600739 & 辽宁成大 & 低位盈利优先；医药生物 / 单日反弹 & H1扣非 15.8亿；Q1 OCF -1.8亿；证据中；概率空间 290.4\% & 补齐当期原始研报锚；验证H2扣非兑现和现金流。 \\
601360 & 三六零 & 低位盈利优先；计算机 / 资金观察 & H1扣非 1.6亿；Q1 OCF -0.6亿；证据中；概率空间 45.6\%；PS on 2026E revenue：2026E收入95.91亿元；熊/基准/牛PS倍数为7.5、10.0、11.0倍，对应熊值10.28元、基准值13.70元、牛值15.07元；概率价值12.95元，条件性目标13.10元。 & 补齐当期原始研报锚；验证H2扣非兑现和现金流。 \\
600918 & 中泰证券 & 低位盈利优先；非银金融 / 低位无流 & H1扣非 17.0亿；Q1 OCF 49.2亿；证据中；概率空间 -6.3\% & 当前不升级；先验证盈利质量、估值消化和现金流反转。 \\
000975 & 山金国际 & 低位盈利优先；有色金属 / 单日反弹 & H1扣非 23.4亿；Q1 OCF 24.9亿；证据中高；概率空间 5.0\% & 需要安全边际；验证H2分母不下修、价格承接有效。 \\
000623 & 吉林敖东 & 低位盈利优先；医药生物 / 单日反弹 & H1扣非 22.2亿；Q1 OCF 0.1亿；证据中；概率空间 317.0\% & 补齐当期原始研报锚；验证H2扣非兑现和现金流。 \\
600292 & 电投水电 & 低位盈利优先；公用事业 / 低位无流 & H1扣非 11.8亿；Q1 OCF 5.2亿；证据中；概率空间 -47.5\% & 当前不升级；先验证盈利质量、估值消化和现金流反转。 \\
000987 & 越秀资本 & 低位盈利优先；非银金融 / 低位无流 & H1扣非 23.3亿；Q1 OCF 39.1亿；证据中；概率空间 2.5\% & 需要安全边际；验证H2分母不下修、价格承接有效。 \\
000155 & 川能动力 & 低位盈利优先；公用事业 / 低位无流 & H1扣非 6.8亿；Q1 OCF -3.2亿；证据高；概率空间 -9.5\% & 正式模型：H2扣非、现金流与外部锚同时复核。 \\
601377 & 兴业证券 & 低位盈利优先；非银金融 / 低位无流 & H1扣非 20.9亿；Q1 OCF 105.9亿；证据中高；概率空间 -0.8\% & 当前不升级；先验证盈利质量、估值消化和现金流反转。 \\
600120 & 浙江东方 & 低位盈利优先；非银金融 / 低位无流 & H1扣非 8.4亿；Q1 OCF 0.6亿；证据中低；概率空间 2.2\% & 需要安全边际；验证H2分母不下修、价格承接有效。 \\
000703 & 恒逸石化 & 启动后仍有空间；石油石化 / 启动确认 & H1扣非 57.2亿；Q1 OCF -25.7亿；证据高；概率空间 61.7\% & 正式模型：H2扣非、现金流与外部锚同时复核。 \\
000301 & 东方盛虹 & 启动后仍有空间；石油石化 / 启动确认 & H1扣非 44.2亿；Q1 OCF 35.3亿；证据中高；概率空间 -3.7\% & 当前不升级；先验证盈利质量、估值消化和现金流反转。 \\
001339 & 智微智能 & 启动后仍有空间；计算机 / 资金观察 & H1扣非 3.7亿；Q1 OCF -5.3亿；证据中高；概率空间 -26.7\% & 当前不升级；先验证盈利质量、估值消化和现金流反转。 \\
002558 & 巨人网络 & 启动后仍有空间；传媒 / 资金观察 & H1扣非 23.0亿；Q1 OCF 13.2亿；证据中高；概率空间 15.0\% & 需要安全边际；验证H2分母不下修、价格承接有效。 \\
\bottomrule
\end{longtable}
\normalsize

其中，正向空间但未进入正式模型的标的，已经给出House条件性目标或明确的上市/证据边界；未进入正式模型的原因是外部锚、研报时效、证据质量或现金流闭环尚未达到正式审计门槛，不再以空白单元格表示。
